\section*{Purpose and Organization of the Forward Models}

The estimation methods developed later require a map from OD-time demand to
the counts that would be observed under that demand. This part constructs that
map in stages. Once routing is fixed, the computational objective is to obtain
and reuse the measurement-response matrix $B$, or an operator that applies it
without reconstructing routes. The detailed assignment model is the behavioral reference: it
represents scheduled events, access, transfers, generalized costs, route
choice, and passenger-flow propagation on the complete time-dependent network.
The fixed-routing response model expresses the same estimation interface in
terms of scheduled events and their expected measurement responses. The
implementation-specific timetable-journey construction is outside the public
forward-model contract; this report therefore states the response assumptions
without prescribing a routing backend.

The remaining sections describe two complementary layers. The additive
operator decomposition exploits linearity once routing is fixed, allowing the
same response to be evaluated in complete, compressed, or sharded form without
changing its mathematical meaning. The observation model then maps the
predicted events to count means and supplies the statistical likelihood. Thus
the detailed and reduced assignments are alternative representations of the
routing response, while operator decomposition and the likelihood are
downstream components that can be combined with either representation when
their assumptions match. They are not competing demand models; those are
introduced separately in Part III.
